\documentclass[aspectratio=169,12pt]{beamer}
\usetheme{Madrid}
\usecolortheme{whale}
\usepackage[T1]{fontenc}
\usepackage[utf8]{inputenc}
\usepackage[brazil]{babel}
\usepackage{graphicx}
\usepackage{booktabs}
\usepackage{xcolor}

\definecolor{labrec}{RGB}{20,60,110}
\setbeamercolor{palette primary}{bg=labrec,fg=white}
\setbeamercolor{palette secondary}{bg=labrec!80,fg=white}
\setbeamercolor{palette tertiary}{bg=labrec!60,fg=white}
\setbeamercolor{title}{bg=labrec,fg=white}
\setbeamercolor{frametitle}{bg=labrec!90,fg=white}
\setbeamercolor{structure}{fg=labrec}
\setbeamerfont{title}{size=\Large,series=\bfseries}
\setbeamerfont{subtitle}{size=\normalsize}

\setbeamertemplate{navigation symbols}{}
\setbeamertemplate{footline}{%
  \leavevmode\hbox{%
    \begin{beamercolorbox}[wd=.33\paperwidth,ht=2.5ex,dp=1.5ex,center]{palette primary}%
      \usebeamerfont{author in head/foot}\insertshortauthor
    \end{beamercolorbox}%
    \begin{beamercolorbox}[wd=.34\paperwidth,ht=2.5ex,dp=1.5ex,center]{palette secondary}%
      \usebeamerfont{title in head/foot}\insertshorttitle
    \end{beamercolorbox}%
    \begin{beamercolorbox}[wd=.33\paperwidth,ht=2.5ex,dp=1.5ex,right]{palette tertiary}%
      \insertframenumber{} / \inserttotalframenumber\hspace*{2ex}
    \end{beamercolorbox}%
  }\vskip0pt%
}

\title{OCRSD}
\subtitle{Only Civil Registry Synthetic Dataset\\[4pt]
{\small Geração de documentos sintéticos do registro civil brasileiro}}
\author{LabREC\,---\,Laboratório de Reconhecimento e Extração de Conhecimento}
\institute{UFSC\,---\,Universidade Federal de Santa Catarina}
\date{Outubro de 2026}
\titlegraphic{%
  \vspace{-8pt}%
  \includegraphics[height=1.1cm]{ocrsd-logo.pdf}\hspace{12pt}%
  \includegraphics[height=1.1cm]{labrec-logo.pdf}\hspace{12pt}%
  \includegraphics[height=1.1cm]{registro-civil-logo.pdf}%
}

\begin{document}

% ---- SLIDE 1: Título ----
\begin{frame}[plain]
  \titlepage
\end{frame}

% ---- SLIDE 2: O que é o OCRSD ----
\begin{frame}[shrink=5]{O que é o OCRSD}
\begin{columns}[T]
\begin{column}{0.55\textwidth}
\textbf{Problema:} modelos de OCR e extração de dados precisam de milhares de
documentos reais para treinamento, mas documentos do registro civil contêm
dados pessoais protegidos.

\vspace{8pt}
\textbf{Solução:} o OCRSD gera documentos sintéticos com dados fictícios e
coerentes, escritos sobre os formulários oficiais.

\vspace{8pt}
\begin{itemize}
  \item Modelos padrão: \textbf{DNV} (69 campos) e \textbf{DO} (103 campos)
  \item Registro semântico com nomes, datas, endereços e vínculos consistentes
  \item Escrita por fonte ou modelos neurais locais
  \item Ground truth por campo com coordenadas e valores
  \item Geração determinística (mesma semente = mesmo resultado)
\end{itemize}
\end{column}
\begin{column}{0.42\textwidth}
  \includegraphics[width=\textwidth]{dnv_trad.png}\\[2pt]
  {\tiny Exemplo: DNV sintética (geração tradicional)}
\end{column}
\end{columns}
\end{frame}

% ---- SLIDE 3: Geração assistida por IA ----
\begin{frame}[shrink=5]{Geração assistida por IA local}
\begin{columns}[T]
\begin{column}{0.55\textwidth}
Um modelo de linguagem local (\textbf{qwen3:1.7b} via Ollama, GPU) propõe
registros coerentes que passam pelas mesmas regras de validação do motor
tradicional.

\vspace{6pt}
\textbf{O LLM não gera o PDF.} Ele propõe apenas valores semânticos (nomes,
idades, municípios) que são validados, aceitos ou rejeitados, e depois
escritos pelo motor existente.

\vspace{6pt}
\begin{tabular}{@{}lcc@{}}
\toprule
                        & \textbf{DNV} & \textbf{DO} \\
\midrule
Campos com IA           & 9 de 11      & 8 de 9      \\
Aceitos na 1ª tentativa & 10/10        & 10/10       \\
Tempo (10 docs, GPU)    & 109\,s       & 85\,s       \\
\bottomrule
\end{tabular}

\vspace{6pt}
\begin{itemize}
  \item Proveniência completa (modelo, digest, cenário, tentativas)
  \item Reprodução exata sem chamar a IA de novo
  \item Funciona sem API externa (Ollama local, GTX 1060)
\end{itemize}
\end{column}
\begin{column}{0.42\textwidth}
  \includegraphics[width=\textwidth]{do_ai.png}\\[2pt]
  {\tiny Exemplo: DO sintética (geração assistida por IA)}
\end{column}
\end{columns}
\end{frame}

% ---- SLIDE 4: Arquitetura e interface ----
\begin{frame}[shrink=10]{Arquitetura e interface web}
\begin{columns}[T]
\begin{column}{0.52\textwidth}
\small
\textbf{Plataforma web} (FastAPI + HTMX, sem Node):
worker com fila SQLite e progresso em tempo real;
editor de mapeamentos com assistente de IA;
galeria de modelos, presets de escrita, compartilhamento.

\vspace{6pt}
\textbf{Motor semântico:} entidades com geradores, vínculos e
restrições; mapeamento por inferência ou vínculo explícito.

\vspace{6pt}
\textbf{Demonstração:} 40 documentos
(DNV + DO $\times$ tradicional + IA)
com ground truth, manifesto e proveniência.

\vspace{6pt}
\footnotesize
\begin{tabular}{@{}ll@{}}
\toprule
\textbf{Componente} & \textbf{Tecnologia} \\
\midrule
Interface / Geração & FastAPI, PyMuPDF \\
Escrita & Fonte / TF / PyTorch \\
IA local & Ollama (qwen3:1.7b) \\
Dados / Implantação & SQLite, Docker \\
\bottomrule
\end{tabular}
\end{column}
\begin{column}{0.45\textwidth}
\centering
\vspace{-4pt}
\includegraphics[width=0.48\textwidth]{dnv_ai.png}%
\hfill%
\includegraphics[width=0.48\textwidth]{do_trad.png}\\[2pt]
{\tiny DNV (IA) \hspace{1.5cm} DO (tradicional)}
\end{column}
\end{columns}
\end{frame}

\end{document}
